% ============================================================
% CHAPTER 7 -- CONCLUSION
% FEMCARE: An AI-Powered Women's Reproductive Health Platform
% ============================================================

\chapter{Conclusion}\doublespacing
\label{chap:conclusion}

This project addressed the need for an integrated digital platform that supports women in monitoring their reproductive health, understanding associated health conditions, and accessing local healthcare services. FEMCARE was designed and implemented as a full-stack web application that brings these functions together in a single authenticated system.

The platform integrates menstrual cycle tracking with real-time Dashboard visualisation, an 18-feature PCOS risk assessment using a pre-trained XGBoost classifier, and an AI health assistant that delivers topic-restricted, context-aware responses using the Groq LLM API. A Find Care module enables users to discover local women's healthcare providers, book appointments from daily time slots, and receive automated email confirmations. A Women's Health Awareness module provides searchable educational content on eight reproductive health conditions. All user data is stored in a Supabase-managed PostgreSQL database with Row Level Security enforced at the database level. Functional testing confirmed that all specified requirements operate correctly in the running application.

FEMCARE demonstrates the practical integration of machine learning, large language model assistance, cloud database management, and modern full-stack web development in a women's health application. The completed system provides a cohesive platform through which users can track their menstrual health, complete a structured PCOS assessment, obtain educational health guidance, and book consultations with local providers --- all within a single secure application. The project successfully fulfils the objectives established at the outset and delivers a functional, integrated solution to the problem it set out to address.
